\section{C14.4: extend the unmatched-Bob certificate to the entire small-$e$ source}
\label{sec:recovery-c144-small-e}

This is a partial-domain theorem. It retains the original physical C14.4
source, all strict gaps, nonnegative event weights, the original total-scale
and prior-complement intersections, and independent source and testing fields.
Its source has $x,L>0$ infinitesimal, $\zeta=L/x$ noninfinitesimal and
$\eta$ unbounded. Write $\sigma=1/\eta$, $S=\sigma+L$ and
\[
 e=\frac{xh}{\sigma},\quad E=1-e,\quad
 \kappa=\frac{x}{S},\quad \nu=\frac{\sigma}{S},\quad
 \ell=\frac{L}{S},\quad \nu+\ell=1.
 \tag{SE1}
\]
The letter $\kappa$ in this section denotes $x/S$, not the different
coefficient denoted by $\kappa$ in the R19 reciprocal-$L$ chart.
All source inequalities are taken in the coefficient field. Real residues
are compared only after taking residues separately in the two fields.

\paragraph{The exact source input.}
Use ES1--ES10 from the retained R13 source proof. In particular
$1/18\le E<1$, $0<h,r<1$, $\theta=1-h$, and
\[
\begin{gathered}
 J=1+xh-\sigma(1-r),\qquad M=1+xr-\theta/\zeta>0,\\
 Q=1-Lr,\quad H_0=1+L(1-r),\quad B=P+x(J-\sigma),\\
 P=\sigma(1+L+2xh+rxL+rx^2h)+(1+x)L+x^2h\theta,\\
 \mathcal H=(1+xr)H_0+(1-r)xh(2+xr)
       -e\{2+x(1+h+r\theta)\}+(1+x)e^2.
\end{gathered} \tag{SE2}
\]
The opening-safe sphere identities are
\[
\begin{gathered}
 2\theta P\beta=hB-\sigma H_0,\qquad
 d_0^2+b_0^2=1,\quad d_0,b_0\ge0,\\
 D=\frac{\sigma^2\mathcal H}{2PJ}d_0^2,\qquad
 x\theta\beta=\frac{\sigma^2\mathcal H}{2P}b_0^2,\\
 e\frac BS=\kappa H_0+\nu\mathcal H b_0^2,\qquad
 e^2\frac BS=\kappa(eH_0+h\mathcal H b_0^2).
\end{gathered}\tag{SE3}
\]
Here $P/S\sim1$, $J,Q,H_0\sim1$, $B/S-(1+\kappa)$ is
infinitesimal, and $\mathcal H-E^2$ is infinitesimal. In particular
$B/S$ and $\mathcal H$ are positive bounded units. These estimates
are the already admitted ES bounds; they do not use bounded $\eta$.

\paragraph{Source dichotomy.}
On the entire C14.4 source,
\[
 e\text{ is infinitesimal}\quad\Longleftrightarrow\quad
 \kappa=x/S\text{ is infinitesimal}.\tag{SE4}
\]
Indeed the first nonnegative sum in SE3 gives $\kappa=O(e)$,
and the second gives $e^2=O(\kappa)$. Both implications are
ordered-field statements entirely inside the source field.

Suppose $e$ is infinitesimal. If $\zeta$ is unbounded, then
$\theta/\zeta$ is infinitesimal and $M\sim1$ directly. If $\zeta$
is a positive bounded unit, then the exact relation
\[
                  \ell=\zeta\kappa\tag{SE5}
\]
gives $\ell$ infinitesimal and $\nu\sim1$. Consequently
$B/\sigma\sim1$. Since $\beta\ge0$, the first identity of SE3
implies $hB\ge\sigma H_0$, whence $h\sim1$ and $\theta$ is
infinitesimal. Thus $M\sim1$ in this case as well. We have proved,
uniformly on the whole small-$e$ source,
\[
                    E\sim1,\quad M\sim1,\quad m\in\operatorname{Inf}_{>0}.
                    \tag{SE6}
\]
For the last statement use the exact original coefficient
\[
 m=\frac{(1+x)(1+\sigma r)eS}{2M\zeta P};\tag{SE7}
\]
$1/\zeta$ is bounded and every other coefficient apart from $e$ is
bounded, including the reciprocal of $M$ just proved to be a unit.
No assumption on the ratio of a source infinitesimal to a testing
infinitesimal occurs.

The two nonnegative terms in the first sum in SE3 imply
$\nu b_0^2$ infinitesimal. Hence the exact sphere formula and the
exact formula for $\alpha$ yield
\[
 \frac DS-\frac{\nu^2}{2}\in\operatorname{Inf},\qquad
 \frac\alpha S-\frac{\ell^2}{2}\in\operatorname{Inf},\qquad
 SC(1-r)\sim\tfrac12,\qquad SAr\sim\tfrac12.\tag{SE8}
\]
For clarity the second formula is obtained without dividing by $\ell$:
\[
 \frac\alpha S=
 \frac{(1+\sigma r)\{\ell(1+xr)e+\ell^2ME\}}
      {2Q(P/S)}.\tag{SE9}
\]
These are additive bounded residue estimates, including $D=0$ whenever
that boundary has points in the stated domain. They do not assert a
relative estimate at a vanishing pivot.

\paragraph{Theorem: every unmatched Bob residue on the small-$e$ source.}
Let $e$ be infinitesimal. In the testing field choose normalized homogeneous
Bob coordinates $(c,s)$ with $c^2+s^2=1$. For every independent projective
test and every real or complex state satisfying
\[
                         \overline{s^2}\ne\overline\ell,\tag{SE10}
\]
the complete original C14.4 parent is a finite sum of squares in
$K_c\otimes_{\mathbb R}K_t$.

\paragraph{Proof with the complete parent retained.}
Noninfinitesimal Alice tests use the unchanged adjacent-$C,A$ theorem:
SE8 and $S$ infinitesimal imply that both $C$ and $A$ are unbounded.
For infinitesimal or zero $u$, apply the exact EB3--EB7 homogeneous
frame from the retained R18 proof. Those formulas are algebraic identities;
their only source-limit inputs to the bounded-Gram proof are precisely
$m$ infinitesimal and SE8, now established on the larger source domain.
To make the imported implication explicit, put
\[
 H=S(C+A),\quad \rho=\frac C{C+A},\quad a_h=\frac{SCA}{C+A},
 \qquad p=\frac{uW}{\sqrt S},\ q=\sqrt S Z,\ y=\frac Y{\sqrt S}.
\]
Then $1/\sqrt H$ is bounded, $0<\rho<1$, $a_h\sim1/2$, and
the harmonic normal $N_y=\sqrt H(y-\rho p)$ can be substituted
by $y=\rho p+N_y/\sqrt H$ with bounded source coefficients.
The remaining normals are the EB4 $N_m,N_b,N_r$, with only pure-field
radicals and the positive testing denominators $1+u^2$ and $1+u^2s^2$.
The exact full quarter-parent is EB7,
\[
 N_y^2+a_hp^2+ASu^2y^2+M_B^2+D_B^2+A_B^2+\beta(X+uZ)^2
                         +\tfrac12 K_B.\tag{SE11}
\]
All EB5 core terms and all EB6 event amplitudes are retained. In particular
no state graph is imposed and neither first-normal coupling is removed.
The two positive rows $ASu^2y^2$ and $\beta(X+uZ)^2$ are left whole.
On the other six roots $(N_y,N_m,N_b,N_r,p,q)$ the bounded Gram has
residue equal to $\operatorname{diag}(1,1,1/2,1/2)$ plus the scalar block
\[
 \frac12\begin{pmatrix}
 1&-\bar s\bar c\\
 -\bar s\bar c&\bar\nu^{\,2}\bar s^{\,2}
                            +\bar\ell^{\,2}\bar c^{\,2}
 \end{pmatrix}.
\tag{SE12}
\]
Its determinant is $(\overline{s^2}-\bar\ell)^2/4$, so SE10
makes this finite coefficient Gram strictly positive definite. The
separated bounded-Gram lemma supplies tensor SOS; restoring the two
whole rows gives the complete parent. The exact homogeneous formulas
allow $c=0$, $s=0$ and $u=0$. Signs and real/imaginary state parts are
handled exactly as in EB, and the original positive source and testing
normalizations are pure-field factors. This proves the stated domain.

\paragraph{New part, and exact retained complement.}
Unbounded $\zeta$ already forces $\kappa\le1/\zeta$ infinitesimal,
so the new theorem contains the entire R18 unmatched-residue admission.
Its additional domain is $\zeta$ a positive bounded unit, $e$
infinitesimal, and Bob sine square with nonzero residue; this includes
Bob infinity. In this added domain $\ell$ has zero residue, so the
remaining Bob chart is exactly finite infinitesimal $v$ (including zero,
which may be handled by independent literal-axis exits).

After all these admissions, the C14.4 remainder is contained in the union
of (i) $e$ infinitesimal and matched Bob residue, and (ii) $\zeta$ a
positive bounded unit and $e$ a positive bounded unit. All original
intersections remain in both. In case (ii), SE7 shows that $m$ has a
positive real lower bound, while its upper bound may diverge through
$M=0$. This is a source estimate, not a full-parent certificate. Neither
remaining domain is declared empty or closed. C14.4 and U1 remain open.

\paragraph{Verification scope.}
The new checker verifies the source rescalings and exact identities used
above. The seven EB identities are replayed separately from an unchanged
temporary copy. Their algebra does not by itself prove the source bounds,
strictness, or universal scope; those are the arguments just given.
